% !PRINTED 107
\weberpage{107}

\begin{center}
{\large \S.\ 32.}

\vspace{0.5em}
{\large\gesperrt{Change of Sign of $f(x)$. Roots of Equations of Odd Degree and of Pure Equations.}}
\end{center}

On the basis of the theorems of the preceding paragraph we can prove the following important theorem.

\medskip
\noindent VIII. \gesperrt{If the coefficients of $f(x)$ are real and there exist two real values $a,b$ of $x$ such that $f(a)$ and $f(b)$ have opposite signs, then between $a$ and $b$ there is at least one root $\xi$ of the equation $f(x)=0$.}
\medskip

We shall assume that $a<b$ and that $f(a)$ is negative, $f(b)$ positive. We divide the numbers between $a$ and $b$ into two parts $\mathfrak A$ and $\mathfrak B$ as follows: a number $\alpha$ belongs to $\mathfrak A$ if, between $\alpha$ and $a$, including the endpoints, there is no $x$ for which $f(x)$ becomes positive; and a number $\beta$ belongs to $\mathfrak B$ if between $\beta$ and $a$ there is at least one value of $x$ for which $f(x)$ becomes positive. A third possibility is plainly not possible, and every number between $a$ and $b$ is therefore placed in $\mathfrak A$ or in $\mathfrak B$. We also count the numbers smaller than $a$ as belonging to $\mathfrak A$, and those greater than $b$ as belonging to $\mathfrak B$.

If $\alpha$ belongs to $\mathfrak A$, then every $x$ smaller than $\alpha$ also belongs to $\mathfrak A$; and if $\beta$ belongs to $\mathfrak B$, every larger $x$ belongs to $\mathfrak B$. Thus every $\beta$ is greater than every $\alpha$, and the two numerical domains $\mathfrak A,\mathfrak B$ form a \gesperrt{cut} (Introduction, p. 5), generated by a number $\xi$, such that every number $\alpha$ smaller than $\xi$ belongs to $\mathfrak A$, and every number $\beta$ greater than $\xi$ belongs to $\mathfrak B$.

If now, for some number $x_0$, the value of $f(x_0)$ is negative or positive, then by V of the preceding paragraph a positive magnitude $\varepsilon$ can be so determined that for every $x$ between $x_0-\varepsilon$ and $x_0+\varepsilon$ the values of $f(x)$ are always negative or always positive. It follows from this that $f(\xi)$ can be neither negative nor positive and must therefore be zero. For if $f(\xi)$ were negative, then $f(x)$ would also be negative between $\xi$ and $\xi+\varepsilon$; these values would therefore belong to $\mathfrak A$, although they are greater than $\xi$. If $f(\xi)$ were positive, then $f(x)$ would be positive as long as $x$ lay between $\xi-\varepsilon$ and $\xi$; these values, although smaller than $\xi$, would therefore belong to $\mathfrak B$. Both alternatives are impossible by the definition of
%CONT>>
